\begin{pastexam}[label={px:df3-103-ds-05}]{地方三等資訊處理 103 年 第五題}
（申論；20 分）

給定下列尚未排序之數列：80, 24, 11, 47, 19, 91, 2, 32, 85, 7, 16, 36, 99, 52, 41，請以泡沫排序法（bubble sort）
及快速排序法（quick sort）分別將該數列由小到大排序：
\begin{itemize}[leftmargin=2.2em, nosep]
\item[㈠] 請依序寫出泡沫排序法前五回合的排序結果。（10 分）
\item[㈡] 請依序寫出快速排序法前五回合的排序結果，每一回合用一個樞紐（pivot），並把每一回合所用的樞紐圈起來。
（10 分）
\end{itemize}
\end{pastexam}

\begin{pastsoln}
\textbf{㈠ 泡沫排序}（資結書第~14.2~節、虛擬碼~14.1）：每一回合由左到右比較相鄰兩數，前面比較大就交換；
第 $i$ 回合結束時，最大的 $i$ 個數已經在最後面（下面用粗體標出）。 共 $15$ 個數，第 $i$ 回合比較 $15 - i$ 次。
\begin{center}
\small
\begin{tabular}{cl}
\toprule
回合 & 結果 \\
\midrule
原始 & 80, 24, 11, 47, 19, 91, 2, 32, 85, 7, 16, 36, 99, 52, 41 \\
1 & 24, 11, 47, 19, 80, 2, 32, 85, 7, 16, 36, 91, 52, 41, \textbf{99} \\
2 & 11, 24, 19, 47, 2, 32, 80, 7, 16, 36, 85, 52, 41, \textbf{91}, \textbf{99} \\
3 & 11, 19, 24, 2, 32, 47, 7, 16, 36, 80, 52, 41, \textbf{85}, \textbf{91}, \textbf{99} \\
4 & 11, 19, 2, 24, 32, 7, 16, 36, 47, 52, 41, \textbf{80}, \textbf{85}, \textbf{91}, \textbf{99} \\
5 & 11, 2, 19, 24, 7, 16, 32, 36, 47, 41, \textbf{52}, \textbf{80}, \textbf{85}, \textbf{91}, \textbf{99} \\
\bottomrule
\end{tabular}
\end{center}
以第 1 回合為例：$80$ 一路和 $24, 11, 47, 19$ 交換往右，遇到 $91$ 停下；接著換 $91$ 往右，和 $2, 32, 85, 7, 16, 36$
交換，遇到 $99$ 停下；最後 $99$ 和 $52, 41$ 交換，沉到最後。

\textbf{㈡ 快速排序。} 題目沒有指定取樞紐和分割的方法，這裡用最常見的寫法（資結書第~14.5~節）：取子陣列的
\textbf{第一個數}當樞紐；$i$ 從左往右找第一個大於樞紐的數，$j$ 從右往左找第一個不大於樞紐的數，$i < j$ 就交換，
直到 $i$、$j$ 交錯；最後把樞紐和 $j$ 位置的數交換，樞紐就到了最終位置。 一回合是一次分割，先處理左邊的子陣列
再處理右邊（長度 $\le 1$ 的子陣列不必分割，不算一回合）。 下面用 \fbox{$\,\cdot\,$} 表示圈起來的樞紐，
$[\ ]$ 標出這一回合處理的範圍。
\begin{itemize}
\item \textbf{第 1 回合}，範圍 $[0..14]$，樞紐 $80$：交換 $91 \leftrightarrow 41$、$85 \leftrightarrow 52$；$i$、$j$ 交錯時
$j$ 停在 $36$，樞紐和 $36$ 交換。\\
$[36, 24, 11, 47, 19, 41, 2, 32, 52, 7, 16, \fbox{80}, 99, 85, 91]$
\item \textbf{第 2 回合}，範圍 $[0..10]$，樞紐 $36$：交換 $47 \leftrightarrow 16$、$41 \leftrightarrow 7$；樞紐和 $32$ 交換。\\
$[32, 24, 11, 16, 19, 7, 2, \fbox{36}, 52, 41, 47], 80, 99, 85, 91$
\item \textbf{第 3 回合}，範圍 $[0..6]$，樞紐 $32$：右邊找到的第一個不大於 $32$ 的是 $2$，左邊找不到大於 $32$ 的，
不交換；樞紐和 $2$ 交換。\\
$[2, 24, 11, 16, 19, 7, \fbox{32}], 36, 52, 41, 47, 80, 99, 85, 91$
\item \textbf{第 4 回合}，範圍 $[0..5]$，樞紐 $2$：其他數都比 $2$ 大，$j$ 退到樞紐自己，位置不變。\\
$[\fbox{2}, 24, 11, 16, 19, 7], 32, 36, 52, 41, 47, 80, 99, 85, 91$
\item \textbf{第 5 回合}，範圍 $[1..5]$，樞紐 $24$：其他數都比 $24$ 小，不交換；樞紐和最後的 $7$ 交換。\\
$2, [7, 11, 16, 19, \fbox{24}], 32, 36, 52, 41, 47, 80, 99, 85, 91$
\end{itemize}
五回合後的數列是 $2, 7, 11, 16, 19, 24, 32, 36, 52, 41, 47, 80, 99, 85, 91$；圈起來的 $80, 36, 32, 2, 24$ 都已在最終
位置。 之後還要繼續處理 $[7, 11, 16, 19]$、$[52, 41, 47]$、$[99, 85, 91]$，最後得到
$2, 7, 11, 16, 19, 24, 32, 36, 41, 47, 52, 80, 85, 91, 99$。 用別的分割法（例如 Lomuto，或取最後一個數當樞紐）
中間結果會不同，作答時要先寫明用哪一種。
\end{pastsoln}
